% draft_introduction.tex
% Style: Academic Paper
% Target: 2--3 pages (~1,400--1,800 words of body text)

\section{Introduction}
\label{sec:introduction}

Does social media sentiment causally affect intraday stock price movements? This paper provides
evidence on that question using a daily Bloomberg panel of S\&P~500 firms over 2024--2025,
employing two identification strategies---firm fixed-effects OLS and Double Machine Learning
(DoubleML)---to estimate the effect of Twitter/X sentiment on open-to-close returns.

The question matters for two reasons. First, the structure of retail equity participation has
changed. The share of U.S. equity volume attributable to retail traders roughly doubled between
2010 and 2020 \citep{welch2022}, and platforms such as Twitter/X, Reddit, and StockTwits have
become primary venues through which retail investors form expectations and coordinate trading
activity \citep{cookson2023}. Whether the information transmitted over these platforms is
efficiently incorporated into prices---or whether it induces noise-driven mispricing---has direct
implications for market quality, regulatory oversight, and the design of disclosure rules. Second,
the causal question remains empirically open. A substantial literature documents that
text-derived sentiment correlates with subsequent returns \citep{tetlock2007, bollen2011,
antweiler2004}, but correlation is not causation. Firms that attract more social media attention
may differ systematically from those that do not; sentiment itself may simply respond to price
changes rather than cause them. Resolving this requires identification strategies capable of
isolating the sentiment effect from confounders and from reverse causality.

Prior work has taken steps in this direction. \citet{tetlock2007} shows that the fraction of
negative words in the Wall Street Journal \textit{Abreast of the Market} column predicts
next-day Dow Jones returns and trading volume, with the effect reversing within a week. Building
on this media-sentiment literature, \citet{antweiler2004} find that message board activity on
Yahoo Finance and Raging Bull predicts next-day volatility and, more modestly, returns.
\citet{bollen2011} report that Twitter mood states---extracted from a large corpus of public
tweets---improve the accuracy of next-day DJIA forecasts, with calm and happy states preceding
market advances. More recently, \citet{gu2020} use Fama-MacBeth cross-sectional regressions to
show that lagged Twitter sentiment predicts open-to-open stock returns in a 2015--2017 sample,
with the effect concentrated at a one-day lag consistent with rapid information absorption.
\citet{teti2019} examine the relationship between daily tweet counts, polarity-broken tweet
volumes, and stock-level returns, finding that negative tweet counts predict returns at short
horizons. These studies provide the empirical template for the present analysis.

The gap that motivates this paper is causal identification. The studies cited above rely on
reduced-form regressions---OLS, panel FE, Fama-MacBeth---that account for observable confounders
but are not designed to partial out high-dimensional nonlinear relationships among controls. When
the confounder set includes variables like intraday price range, momentum indicators, capital
structure ratios, and multiple sentiment channels simultaneously, the linearity assumption
embedded in OLS can bias the treatment coefficient in ways that are difficult to diagnose. More
fundamentally, the direction of causation between social media sentiment and stock prices is not
settled: sentiment may drive prices, prices may drive sentiment, or both may respond to the same
latent shock. None of the prior literature addresses this directly with a formal placebo design.

This paper makes two contributions. The first is methodological. By applying the Partially Linear
Regression (PLR) model of \citet{chernozhukov2018}, combined with XGBoost as a nuisance learner,
I obtain a Neyman-orthogonal estimate of the sentiment effect that is asymptotically unconfounded
by high-dimensional covariates, regardless of the functional form of their relationship with the
treatment and outcome. The DoubleML estimator achieves this by cross-fitting: nuisance models for
both $E[\text{twitter\_sent} \mid X]$ and $E[\text{return} \mid X]$ are trained on held-out folds,
so the final coefficient is not contaminated by in-sample overfitting. This is, to my knowledge,
the first application of DoubleML to intraday Twitter sentiment and stock returns.

The second contribution is the systematic documentation of a reverse causality concern. I pair the
primary DoubleML estimates with a firm-level placebo test that regresses yesterday's return on
today's Twitter sentiment. Under the null of no reverse causation, this coefficient should be
indistinguishable from zero. It is not: the estimated coefficient is $+2.348$ ($p < 0.01$),
substantially larger in magnitude than any of the forward-looking estimates. This result is
inconsistent with clean causal transmission from sentiment to prices and raises the possibility
that Bloomberg's proprietary sentiment measure largely reflects the market's own prior-day
behavior rather than an independent information signal. Additional lag-decay and impact-persistence
tests reinforce this interpretation.

The main findings can be summarized as follows. In the firm fixed-effects specifications (four
models, $N \approx 160{,}000$ firm-day observations), \texttt{twitter\_sent} is statistically
insignificant across all treatment constructions, including continuous sentiment, a
negative-clipped measure, and negative tweet counts. News sentiment (\texttt{news\_sent}) enters
with a coefficient of $-1.056$ ($p < 0.01$) and the RSI-30 momentum indicator is significant,
suggesting that within-firm variation in price-relevant news and momentum dominate the Twitter
channel in a linear framework. In the DoubleML specifications, the picture is noisier but more
significant. The PLR estimate for \texttt{twitter\_sent} is $-0.924$ ($p < 0.01$), and the
estimate for \texttt{news\_sent} is $-0.759$ ($p < 0.01$). These estimates survive 20 repetitions
of 5-fold cross-fitting with 1{,}000 bootstrap draws. The negative direction of the Twitter
coefficient is consistent with a ``bad news amplification'' interpretation, but the magnitude is
economically modest and the placebo evidence makes causal attribution fragile.

A partial replication of \citet{gu2020} using Fama-MacBeth regressions on risk-adjusted
(four-factor) returns produces a coefficient on lagged Twitter sentiment of $-0.036$, which is
insignificant and directionally opposite to the original paper's $+0.136^{***}$. A long-short
strategy sorting on daily sentiment deciles generates an annualized Sharpe ratio of 1.66 over 329
trading days---roughly half the Sharpe of 3.17 reported in the original 2015--2017 sample. Taken
together, these comparisons suggest that the predictive content of Twitter sentiment for stock
returns may have declined materially since the period studied by \citet{gu2020}, a possibility
examined through temporal analysis in Section~\ref{sec:results}.

The remainder of the paper proceeds as follows. Section~\ref{sec:literature} reviews the relevant
literature on sentiment and asset pricing, social media as an information channel, and causal
identification in high-frequency panels. Section~\ref{sec:data} describes the Bloomberg panel,
variable construction, and sample properties. Section~\ref{sec:methodology} presents the FE-OLS
and DoubleML specifications with full notation and identifying assumptions.
Section~\ref{sec:results} reports point estimates, robustness checks, and the placebo and
lag-structure tests. Section~\ref{sec:discussion} interprets the central tension between the
DoubleML estimates and the reverse causality evidence, and places the findings in the context of
prior literature. Section~\ref{sec:conclusion} concludes with limitations and directions for
future work, including a custom entity-level sentiment tool that may isolate causally relevant
signal from the Bloomberg composite.


% =========================================================
% MISSING / EXPAND NOTES
% =========================================================
% [MISSING: \bibliography{} entry -- no .bib file exists yet.
%   Required keys: gu2020, teti2019, tetlock2007, bollen2011, antweiler2004,
%   chernozhukov2018, welch2022, cookson2023, baker2006, da2011, chen2014,
%   famaMacbeth1973]
% [EXPAND: If a sample period table or summary statistics figure is added
%   to the Data section, consider adding one sentence here stating the
%   observation window explicitly (currently "2024--2025").]
% [EXPAND: If a temporal analysis subsection finds a structural break,
%   revisit the paragraph on Gu & Kurov decline to add the break date.]
